% !TEX root = conj_sys.tex

\section{Aromaticity \& Aromatic Substitution} % lessons 6-9

  \subsection{Aromatic Compounds}

    \begin{core}{Aromatic Compounds}

      The name "aromatic" originates from these compounds' characteristic scent; however, the term now denotes a strictly structural definition, as covered in \autoref{core:aromaticity_rules}.

      While aromaticity is not limited to the base benzene structure, these notes will primarily refer to benzene-derived compounds such as the following.

      \doublereactionfig{
        \scheme{
          \chemname{
            \chemfig[atom sep=2em]{
              *6(=-=
              (-[:30]H)
              -=-)
            }
          }{\textit{Benzene (PhH)}}
          \qquad \qquad
          \chemname{
            \chemfig[atom sep=2em]{
              *6(=-=
              (-[:30]\ce{CH3})
              -=-)
            }
          }{\textit{Toluene (PhMe)}}
          \qquad \qquad
          \chemname{
            \chemfig[atom sep=2em]{
              *6(=-=
              (-[:30]OH)
              -=-)
            }
          }{\textit{Phenol (PhOH)}}
          \qquad \qquad
          \chemname{
            \chemfig[atom sep=2em]{
              *6(=-=
              (-[:30]\ce{NH2})
              -=-)
            }
          }{\textit{Aniline (\ce{PhNH2})}}
        }
      }{
        \scheme{
          \chemname{
            \chemfig[atom sep=2em]{
              *6(=-=
              (
              -[:30]
              =[:-30]O
              )
              -=-)
            }
          }{\textit{Aniline (\ce{PhNH2})}}
          \qquad \qquad
          \chemname{
            \chemfig[atom sep=2em]{
              *6(=-=
              (
              -[:30]
              (=[:90]O)
              -[:-30]OH
              )
              -=-)
            }
          }{\textit{Benzoic Acid (\ce{PhCO2H})}}
          \qquad \qquad
          \chemname{
            \chemfig[atom sep=2em]{
              *6(=-=-N=-)
            }
          }{\textit{Pyridine}}
        }
      }{Phenyl ring-based aromatic compound examples; \textit{pyridine represents an aromatic compound with a non-CH atom as part of its ring}}{fig:ph57a}{0.9}

    \end{core}

    \begin{core}{Aromatic Compound Stereochemistry}

      Additional substituent positions on the root aromatic compounds shown above can be specified using the ortho, meta, and para keywords.

      \reactionfig{
        \scheme{
          \chemfig[atom sep=2em]{
            *6(
            (-[:-150]R)
            =
            (-[:-90,,,,Yellow]\textcolor{Yellow}{ortho})
            -
            (-[:-30,,,,Orange]\textcolor{Orange}{meta})
            =
            (-[:30,,,,Blue]\textcolor{Blue}{para})
            -
            (-[:90,,,,Orange]\textcolor{Orange}{meta})
            =
            (-[:150,,,,Yellow]\textcolor{Yellow}{ortho})
            -)
          }
        }
      }{Phenyl ring substituent-positioning keywords}{fig:ph58a}{0.9}

      The compound naming scheme follows the structure \textit{relation-substituent-root}. Consider the following examples.

      \reactionfig{
        \scheme{
          \chemname{
            \chemfig[atom sep=2em]{
              *6(=
              (-[:-90]\ce{NH2})
              -
              (-[:-30]\ce{NO2})
              =-=-)
            }
          }{\textcolor{Yellow}{ortho-}\textcolor{Blue}{nitro-}\textcolor{Orange}{aniline}}
          \qquad \qquad
          \chemname{
            \chemfig[atom sep=2em]{
              *6(=
              (-[:-90]Cl)
              -=
              (
              -[:30]
              (=[:90]O)
              -[:-30]OH
              )
              -=-)
            }
          }{\textcolor{Yellow}{meta-}\textcolor{Blue}{chloro-}\textcolor{Orange}{benzoic acid}}
          \qquad \qquad
          \chemname{
            \chemfig[atom sep=2em]{
              *6(=-
              (-[:-30]OH)
              =
              (-[:30]OH)
              -=-)
            }
          }{\textcolor{Yellow}{ortho-}\textcolor{Blue}{hydroxyl-}\textcolor{Orange}{phenol}}
          \qquad \qquad
          \chemname{
            \chemfig[atom sep=2em]{
              *6(=-
              (
              -[:-30]N
              (-[:-90])
              -[:30]
              )
              =-=-)
            }
          }{\textcolor{Yellow}{para-}\textcolor{Blue}{dimethylamino-}\textcolor{Orange}{pyridine}}
        }
      }{Phenyl ring-based aromatic compound nomenclature}{fig:ph59a}{0.9}

    \end{core}

    \begin{core}{Reactivity of Aromatic Compounds}

      Aromatic compounds are considerably less reactive than would otherwise be predicted. For instance, while a simple cyclohexene reacts readily with mCPBA to form an epoxide, no reaction occurs when benzene is mixed with three equivalents of mCPBA.

      \doublereactionfig{
        \scheme{
          \chemfig[atom sep=1.5em]{
            *6(--=---)
          }
          \+
          \chemfig{
            \text{mCPBA}
          }
          \arrow{->}
          \chemfig[atom sep=1.5em]{
            *6(--
            (*3(-O--))
            ----)
          }
          \+
          \chemfig{
            \text{mCBA}
          }
        }
      }{
        \scheme{
          \chemfig[atom sep=1.5em]{
            *6(=-=-=-)
          }
          \+
          \chemfig{
            \text{3 mCPBA}
          }
          \arrow{->}
          \chemfig{
            \text{\textit{no reaction}}
          }
        }
      }{Reactions of cyclohexene and benzene with mCPBA}{fig:ph60a}{0.9}
      \vspace{-0.4cm}

      Similarly, a substantial portion of the predicted energy release during the hydrogenation of benzene is offset by aromatic stabilization.

      \pgfig{
        % Baseline (x-axis)
        \draw[thick] (0,0) -- (12,0);
        \draw[-{Latex}, thick] (0,0) -- (0,6);
        \node[rotate=0] at (-0.5,3) {\large E};

        % Molecule structures at top of each column
        \node at (1,3) {
          \chemfig[atom sep=1.5em]{
            *6(--=---)
          }
        };
        \node at (4,5) {
          \chemfig[atom sep=1.5em]{
            *6(--=--=)
          }
        };
        \node at (7,5) {
          \chemfig[atom sep=1.5em]{
            *6(--=-=-)
          }
        };
        \node at (10,6) {
          \chemfig[atom sep=1.5em]{
            *6(=-=-=-)
          }
        };

        % Column 1
        \draw[Yellow, thick, {Bar}->] (1,1.5) -- (1,0) node[midway, right] {-28.6};

        % Column 2
        \draw[Blue, thick, {Bar}->] (3,3) -- (3,0) node[midway, right] {-57.2} node[above] at (3,3) {\footnotesize predicted};

        % Column 3
        \draw[Yellow, thick, {Bar}->] (5,3) -- (5,0) node[midway, right] {-57.2} node[above] at (5,3) {\footnotesize actual};

        % Column 4
        \draw[Yellow, thick, {Bar}->] (7,2.7) -- (7,0) node[midway, right] {-55.4};

        \draw[Orange, thick, {Bar}-{Bar}] (7,3) -- (7,2.7) node[above, align=center] at (7,3) {\footnotesize ~$2\,\text{kcal}\,\text{mol}^{-1}$ \\\footnotesize of conj. stabilization};

        % Column 5
        \draw[Blue, thick, {Bar}->] (9,4.5) -- (9,0) node[midway, right] {-80} node[above] at (9,4.5) {\footnotesize predicted};

        % Column 6
        \draw[Yellow, thick, {Bar}->] (11,2.3) -- (11,0) node[midway, right] {-50};

        \draw[Orange, thick, {Bar}-{Bar}] (11,4.5) -- (11,2.3) node[midway, right, align=left] {\footnotesize $30\,\text{kcal}\,\text{mol}^{-1}$\\\footnotesize of aromatic stabilization};
      }{Energy diagram of predicted versus actual heats of hydrogenation for cyclohexene derivatives; \textit{aromaticity confers a substantial increase in stability relative to the summed values of the corresponding isolated double-bond hydrogenations, exceeding even the stabilization observed in nonaromatic but conjugated cyclohexene systems}}{fig:ph61a}{0.9}
      \vspace{-0.4cm}

      To explain this phenomenon, consider a top-down view of the benzene π molecular orbital diagram.

      \myfig{0.6}{conj_sys_figures/benzene_frost}{Energy diagram of the benzene π molecular orbitals (\textit{top-down view})}{fig:ph62a}

      As shown in \autoref{fig:ph62a}, the π electrons occupy exclusively bonding molecular orbitals, rendering the system exceptionally stable.

    \end{core}

    \begin{core}{Aromatic Compound Molecular Diagrams}

      The \textbf{Frost Diagram} is a method for rapidly predicting π molecular orbital energy levels of a cyclic conjugated system, following the steps below.

      \begin{enumerate}[label=(\arabic*)]

        \item The compound must possess a cyclic array of p orbitals

        \item Draw the ring structure with one vertex pointed downwards

        \item Draw a circle around the compound

        \item Intersections represent energy levels

        \item The middle of the circle represents the nonbonding energy level

        \item Add \ce{e^-}s

      \end{enumerate}

      Below are the corresponding orbital diagrams for a cyclopentadienyl anion and a cycloheptatrienyl cation; the benzene diagram has already been shown above.

      \begin{enumerate}

        \item \textbf{Cyclopentadienyl Anion}

        \pgfig{
          % Cyclopentadienyl anion structure (left)
          \node at (0.3,1.8) {
            \chemfig{
              [:18]*5(-=
              -[:]\ncharge{90}{}
              -=)
            }
          };

          % Dashed correlation line
          \draw[dashed] (1.8,1.8) -- (10.5,1.8);

          % Hückel pentagon energy-level diagram (yellow/orange pentagon outline representing MO pattern)
          \node at (3.2,1.8) {
            \chemfig{
              [:198]*5(-=
              -[:]\ncharge{-90}{}
              -=)
            }
          };
          \draw[Yellow, thick] (3.2,1.9) circle (0.9cm);

          \draw[Orange, thick] (2.9,1) -- (3.5,1);
          \draw[Orange, thick] (2.1,1.6) -- (2.7,1.6);
          \draw[Orange, thick] (3.7,1.6) -- (4.3,1.6);
          \draw[Orange, thick] (2.3,2.6) -- (2.9,2.6);
          \draw[Orange, thick] (3.5,2.6) -- (4.1,2.6);

          % Energy axis
          \draw[thick, ->] (6,0.5) -- (6,3.5) coordinate (Etop) node[anchor=east, xshift=-0.1cm, yshift=-0.4cm] {$E$};

          \begin{scope}[xshift=5cm]
            \draw[Orange, thick] (2.9,1) -- (3.5,1) node[midway, above, Blue] {$\uparrow\,\downarrow$};
            \draw[Orange, thick] (2.1,1.6) -- (2.7,1.6) node[midway, above, Blue] {$\uparrow\,\downarrow$};
            \draw[Orange, thick] (3.7,1.6) -- (4.3,1.6) node[midway, above, Blue] {$\uparrow\,\downarrow$};
            \draw[Orange, thick] (2.3,2.6) -- (2.9,2.6);
            \draw[Orange, thick] (3.5,2.6) -- (4.1,2.6);
          \end{scope}
        }{Frost diagram for the cyclopentadienyl anion, showing three filled bonding/nonbonding orbitals and two unoccupied, degenerate antibonding orbitals}{fig:ph63a}{1}

        \item \textbf{Cycloheptatrienyl Cation}

        \pgfig{
          \node at (0,1.8) {
            \chemfig{
              [:13]*7(
              -[:]\pcharge{-90}{}
              -=-=-=)
            }
          };

          \draw[dashed] (1.8,1.8) -- (10.5,1.8);

          % Hückel pentagon energy-level diagram (yellow/orange pentagon outline representing MO pattern)
          \node at (3.2,1.8) {
            \chemfig{
              [:13]*7(
              -[:]\pcharge{-90}{}
              -=-=-=)
            }
          };
          \draw[Yellow, thick] (3.2,1.85) circle (1.25cm);

          \draw[Orange, thick] (2.9,0.6) -- (3.5,0.6);

          \draw[Orange, thick] (2,1.1) -- (2.6,1.1);
          \draw[Orange, thick] (3.8,1.1) -- (4.4,1.1);

          \draw[Orange, thick] (1.7,2.1) -- (2.3,2.1);
          \draw[Orange, thick] (4.1,2.1) -- (4.7,2.1);

          \draw[Orange, thick] (2.4,3) -- (3,3);
          \draw[Orange, thick] (3.4,3) -- (4,3);

          % Energy axis
          \draw[thick, ->] (6,0.5) -- (6,3.5) coordinate (Etop) node[anchor=east, xshift=-0.1cm, yshift=-0.4cm] {$E$};

          \begin{scope}[xshift=5cm]
            \draw[Orange, thick] (2.9,0.6) -- (3.5,0.6) node[midway, above, Blue] {$\uparrow\,\downarrow$};
            \draw[Orange, thick] (2,1.1) -- (2.6,1.1) node[midway, above, Blue] {$\uparrow\,\downarrow$};
            \draw[Orange, thick] (3.8,1.1) -- (4.4,1.1) node[midway, above, Blue] {$\uparrow\,\downarrow$};

            \draw[Orange, thick] (1.7,2.1) -- (2.3,2.1);
            \draw[Orange, thick] (4.1,2.1) -- (4.7,2.1);

            \draw[Orange, thick] (2.4,3) -- (3,3);
            \draw[Orange, thick] (3.4,3) -- (4,3);
          \end{scope}
        }{Frost diagram for the cycloheptatrienyl cation, showing three filled bonding/nonbonding orbitals and four unoccupied, degenerate antibonding orbitals}{fig:ph64a}{1}

      \end{enumerate}

    \end{core}

    \begin{core}[label=core:aromaticity_rules]{Aromaticity Rules}

      In order for a compound to be considered aromatic, it must satisfy a strict set of rules.

      \begin{enumerate}[label=(\arabic*)]

        \item The compound must be a cyclic array of p orbitals

        \item The \ce{e^-} count of the cyclic array must fulfill the equation of $4n+2$, where $n$ is an integer (0, 1, 2, 3, 4, \ldots).

        \textit{An \ce{e^-} count that instead fulfills the equation $4n$, where $n$ is an integer, results in a compound that is far more reactive than predicted and is termed \textbf{anti-aromatic}.}

      \end{enumerate}

      If the above criteria are not met, the compound is considered \textbf{nonaromatic}, as in the cyclohexene shown below.

      \reactionfig{
        \scheme{
          \chemfig{
            *6(=
            -\charge{-30:4pt=A}{}
            -\charge{30:4pt=B}{}
            -=-)
          }
        }
      }{Nonaromatic cyclohexene; \textit{carbons A and B are sp$^3$, which prevents an uninterrupted cyclic array of p orbitals}}{fig:ph65a}{0.9}

      As practice, consider the following two examples.

      \begin{enumerate}

        \item \textbf{Cyclopentadienyl Anion}

        \reactionfig{
          \scheme{
            \chemfig{
              [:18]*5(-=
              -[:]\ncharge{90}{}
              -=)
            }
          }
        }{Aromatic cyclopentadienyl anion}{fig:ph66a}{0.9}

        This structure exhibits a cyclic array of p orbitals with an electron count of 6, satisfying the $4n+2$ expression. This compound is therefore \textbf{considered aromatic}.

        \item \textbf{Cyclopentadienyl Cation}

        \reactionfig{
          \scheme{
            \chemfig{
              [:18]*5(-=
              -[:]\pcharge{90}{}
              -=)
            }
          }
        }{Anti-aromatic cyclopentadienyl cation}{fig:ph67a}{0.9}

        This compound likewise exhibits a cyclic array of p orbitals, but its electron count of 4 satisfies $4n$ rather than $4n+2$. This compound is therefore considered \textbf{anti-aromatic}.

        Notice that the corresponding molecular orbital diagram results in two unpaired electrons, accounting for its heightened reactivity.

        \pgfig{
          % Cyclopentadienyl anion structure (left)
          \node at (0.3,1.8) {
            \chemfig{
              [:18]*5(-=
              -[:]\pcharge{90}{}
              -=)
            }
          };

          % Dashed correlation line
          \draw[dashed] (1.8,1.8) -- (10.5,1.8);

          % Hückel pentagon energy-level diagram (yellow/orange pentagon outline representing MO pattern)
          \node at (3.2,1.8) {
            \chemfig{
              [:198]*5(-=
              -[:]\pcharge{-90}{}
              -=)
            }
          };
          \draw[Yellow, thick] (3.2,1.9) circle (0.9cm);

          \draw[Orange, thick] (2.9,1) -- (3.5,1);
          \draw[Orange, thick] (2.1,1.6) -- (2.7,1.6);
          \draw[Orange, thick] (3.7,1.6) -- (4.3,1.6);
          \draw[Orange, thick] (2.3,2.6) -- (2.9,2.6);
          \draw[Orange, thick] (3.5,2.6) -- (4.1,2.6);

          % Energy axis
          \draw[thick, ->] (6,0.5) -- (6,3.5) coordinate (Etop) node[anchor=east, xshift=-0.1cm, yshift=-0.4cm] {$E$};

          \begin{scope}[xshift=5cm]
            \draw[Orange, thick] (2.9,1) -- (3.5,1) node[midway, above, Blue] {$\uparrow\,\downarrow$};
            \draw[Orange, thick] (2.1,1.6) -- (2.7,1.6) node[midway, above, Blue] {$\uparrow$};
            \draw[Orange, thick] (3.7,1.6) -- (4.3,1.6) node[midway, above, Blue] {$\uparrow$};
            \draw[Orange, thick] (2.3,2.6) -- (2.9,2.6);
            \draw[Orange, thick] (3.5,2.6) -- (4.1,2.6);
          \end{scope}
        }{Frost diagram for the cyclopentadienyl cation}{fig:ph68a}{1}

      \end{enumerate}

    \end{core}

    \begin{core}{Aromatic Contributors}

      An \textbf{aromatic contributor} refers to a resonance structure of a molecule that possesses a fully conjugated, cyclic, planar π system satisfying the $4n+2$ π-electron count, even when the dominant or lowest-energy Lewis structure does not immediately appear aromatic. Consider the following molecule as an example.

      \reactionfig{
        \scheme{
          \chemfig{
            [:-60]*6(-@{nitrogen1}\charge{-90:12pt=\footnotesize\textcolor{Blue}{p}}{}\charge{-90:4pt=\:}{N}
            (-[:-150]R)
            -[@{bond1}:]
            (=[@{bond2}:-90]@{oxygen1}\charge{-135:16pt=\footnotesize\textcolor{Blue}{sp$^2$}}{}\charge{-135:4pt=\:}{}\charge{-30:20pt=\footnotesize\textcolor{Blue}{sp$^2$}}{}\charge{-45:4pt=\:}{O})
            -@{nitrogen2}\charge{-30:20pt=\footnotesize\textcolor{Blue}{sp$^2$}}{}\charge{-45:4pt=\:}{N}
            =[@{bond3}:]
            (-[@{bond4}:30]@{nitrogen3}\charge{-70:12pt=\footnotesize\textcolor{Blue}{p}}{}\charge{-90:4pt=\:}{\ce{NH2}})
            -=)
          }
          \arrow{<->}
          \chemfig{
            [:-60]*6(-\pcharge{-90}{N}
            (-[:-150]R)
            =
            (-[:-90]\ncharge{-45}{O})
            -N=
            (-[:30]\ce{NH2})
            -=)
          }
          \arrow{->}[-90,0.5,,,White]
        }
        \chemmovearrow{
          \reactionfigarrow{nitrogen1}{60}{bond1}{30}{Red}
          \reactionfigarrow{bond2}{10}{oxygen1}{0}{Red}
        }
      }{Aromatic contributor; \textit{a simple rule of thumb for identifying a resonance contributor is to attempt to achieve alternating π bonds while ensuring that a circle drawn around the ring never crosses a π bond}}{fig:ph69a}{0.85}

      Notice that although the base resonance structure is nonaromatic by definition, since its conjugated π orbital system is not cyclic, an alternative structure accessible through resonance is aromatic, thereby qualifying the molecule as possessing an aromatic contributor.

    \end{core}

    \begin{core}{Aromaticity Exceptions \& Basicity}

      While most aromatic systems conform to the criteria outlined in \autoref{core:aromaticity_rules}, certain ring structures constitute notable exceptions.

      \reactionfig{
        \scheme{
          \chemfig[atom sep=2.5em]{
            *7(-=-=-O-=)
          }
        }
      }{Oxepin; \textit{although its p orbitals appear to form a cyclic array satisfying the $2n$ electron count associated with anti-aromaticity, experimental data show oxepin to be non-aromatic}}{fig:oxepin}{0.9}
      \vspace{-0.4cm}

      An initial assessment of oxepin's (\textit{\autoref{fig:oxepin}}) π system suggests a total of eight p-orbital electrons, six contributed by the three double bonds and two by the oxygen's lone pair, which would nominally classify the molecule as anti-aromatic. However, because anti-aromaticity imposes a severe thermodynamic penalty, the molecule instead avoids this unstable configuration by disrupting the continuous cyclic array required for it. It accomplishes this by \textbf{localizing the oxygen's lone pair and excluding it from resonance}. By breaking the conjugated loop of p orbitals in this manner, oxepin renders itself \textbf{non-aromatic rather than anti-aromatic}.

      \reactionfig{
        \scheme{
          \chemname{
            \chemfig{
              *7(-=-=-O-=)
            }
          }{\textit{no resonance}}
          \qquad
          \text{vs}
          \qquad
          \chemfig{
            *5(
            =
            -@{carbon}
            =[@{bond2}:]
            -[@{bond1}:]@{oxygen}O
            -
            )
          }
          \arrow{<->}
          \chemfig{
            *5(=
            -\ncharge{-90}{}
            -
            =\pcharge{135}{O}
            -)
          }
        }
        \chemmovearrow{
          \reactionfigarrow{oxygen}{60}{bond1}{30}{Red}
          \reactionfigarrow{bond2}{-30}{carbon}{-60}{Red}
        }
      }{Comparison of resonance behavior in oxepin (\textit{left}) versus furan (\textit{right})}{fig:aro_bas}{0.9}

      This structural anomaly carries direct consequences for basicity, as illustrated by the comparison between oxepin and furan (\textit{\autoref{fig:aro_bas}}).

      \begin{itemize}

        \item In \textbf{oxepin}, the oxygen's lone pair is strictly localized in order to prevent the continuous cyclic array that would otherwise produce an unstable anti-aromatic state. Because these electrons remain confined to the oxygen and do not participate in resonance, they stay fully available to accept a proton.

        \item In contrast, the oxygen's lone pair in \textbf{furan} delocalizes actively into the ring's π system. This resonance withdraws electron density from the oxygen, imparting a partial positive charge and substantially diminishing its capacity to act as an electron donor.

      \end{itemize}

      Consequently, the forced localization of the lone pair in oxepin renders its oxygen atom significantly more basic than the resonance-delocalized oxygen of the six-membered furan ring.

    \end{core}

  \subsection{Electrophilic Aromatic Substitution}
    \label{sec:eas}

    \begin{core}[label=core:eas_general]{Electrophilic Aromatic Substitution}

      \textbf{Electrophilic aromatic substitution (EAS)} describes a class of reactions in which an aromatic ring functions as the nucleophilic species. The general mechanism underlying this reaction type is outlined below.

      \reactionfig{
        \scheme{
          \chemfig{
            [:0]E
            -[@{bond1}:0]@{leaving}LG
          }
          \arrow{->}
          \chemfig{
            [:0]@{elec}\pcharge{60}{E}
          }
          \qquad
          \chemfig[atom sep=2em]{
            *6(
            -=-=
            -@{carbon1}
            =[@{bond2}:]
            )
          }
          \arrow{->}
          \chemfig[atom sep=2em]{
            *6(
            (\pcharge{-150}{})
            -=-=-
            (-[:180]E)
            (-[@{bond3}:120]@{hydrogen}H)
            -[@{bond4}:])
          }
          \raisebox{0.8cm}{
            \+
            \chemfig{
              [:0]@{base}Base
            }
          }
          \arrow{->}
          \chemfig[atom sep=2em]{
            *6(-=-=-
            (-[:150]E)
            =)
          }
          \+
          \chemfig{
            [:0]\ncharge{60}{LG}
          }
          \+
          \chemfig{
            [:0]H
            -[:0]\pcharge{90}{Base}
          }
        }
        \chemmovearrow{
          \reactionfigarrow{bond1}{-100}{leaving}{-80}{Red}
          \reactionfigarrow{bond2}{-10}{carbon1}{0}{Red}
          \reactionfigarrow{carbon1}{135}{elec}{100}{Red}
          \reactionfigarrow{base}{135}{hydrogen}{45}{Red}
          \reactionfigarrow{bond3}{30}{bond4}{-15}{Red}
        }
      }{General mechanistic outline of electrophilic aromatic substitution reactions}{fig:ph70a}{0.8}

    \end{core}

    \subsubsection{Reaction Types}
      \label{sec:eas_types}

      \begin{core}[label=core:eas_halogenation]{Halogenation}

        \textbf{EAS halogenation} results in the attachment of a single halogen atom to an aromatic ring, with the corresponding hydrohalic acid generated as a byproduct. This mechanism proceeds via activation by a \textbf{Lewis acid} (\textit{defined as any chemical species capable of accepting an electron pair}), as illustrated below.

        \vspace{0.4cm}
        \reactionfig{
          \scheme{
            \chemfig[atom sep=2.5em]{
              [:0]X
              -[:0]@{halogen1}X
            }
            \qquad
            \chemfig{
              @{la}\text{Lewis Acid}
            }
            \arrow{->}
            \chemfig[atom sep=2.5em]{
              [:0]X
              -[@{bond1}:0]@{halogen2}\pcharge{-90}{X}
              -[:0]\ncharge{90}{LA}
            }
            \arrow{->}
            \chemfig{
              [:0]\pcharge{45}{X}
            }
            \+
            \chemfig[atom sep=2.5em]{
              [:0]X
              -[:0]\ncharge{90}{LA}
            }
          }
          \chemmovearrow{
            \reactionfigarrow{halogen1}{30}{la}{150}{Red}
            \reactionfigarrow{bond1}{100}{halogen2}{80}{Red}
          }
        }{Generation of a halogen cation electrophile via Lewis acid activation; \textit{Lewis acids function as electron-pair acceptors and include species such as \ce{FeX3}, \ce{BX3}, and \ce{AlX3}}}{fig:ph71a}{0.9}
        \vspace{-0.4cm}

        The chlorination and bromination of benzene, presented below, illustrate this general mechanism.

        \begin{enumerate}

          \item \textbf{Chlorination of Benzene}

          \doublereactionfig{
            \scheme{
              \chemfig[atom sep=2.5em]{
                [:0]Cl
                -[:0]@{halogen1}Cl
              }
              \qquad
              \chemfig{
                @{la}\ce{AlCl3}
              }
              \arrow{->}
              \chemfig[atom sep=2.5em]{
                [:0]Cl
                -[@{bond1}:0]@{halogen2}\pcharge{-90}{Cl}
                -[:0]\ncharge{90}{Al}\ce{Cl3}
              }
              \arrow{->}
              \chemfig{
                [:0]\pcharge{45}{Cl}
              }
              \+
              \chemfig[atom sep=2.5em]{
                [:0]Cl
                -[:0]\ncharge{90}{Al}\ce{Cl3}
              }
            }
            \chemmovearrow{
              \reactionfigarrow{halogen1}{-60}{la}{-120}{Red}
              \reactionfigarrow{bond1}{100}{halogen2}{80}{Red}
            }
          }{
            \scheme{
              \chemfig[atom sep=2em]{
                *6(=-=-=-)
              }
              \+
              \chemfig{
                [:0]\ce{Cl2}
              }
              \qquad \qquad
              \arrow{->[cat \ce{AlCl3}][
                \shortstack{
                  electrophilic \\
                  \ce{Cl+} attack
                }
              ]}
              \qquad
              \chemfig[atom sep=2em]{
                *6(
                (\pcharge{-150}{})
                -=-=-
                (-[:180]Cl)
                (-[@{bond3}:120]@{hydrogen}H)
                -[@{bond4}:])
              }
              \+
              \chemfig{
                [:0]@{base}Cl
                -[@{bond5}:0]\ncharge{90}{Al}\ce{Cl3}
              }
              \arrow{->}
              \chemfig[atom sep=2em]{
                *6(-=-=-
                (-[:150]Cl)
                =)
              }
              \+
              \chemfig{
                [:0]HCl
              }
            }
            \chemmovearrow{
              \reactionfigarrow{base}{135}{hydrogen}{45}{Red}
              \reactionfigarrow{bond3}{30}{bond4}{-15}{Red}
              \reactionfigarrow{bond5}{80}{base}{100}{Red}
            }
          }{Chlorination of benzene}{fig:ph72a}{0.9}
          \vspace{-0.4cm}

          \item \textbf{Bromination of Benzene}

          \doublereactionfig{
            \scheme{
              \chemfig[atom sep=2.5em]{
                [:0]Br
                -[:0]@{halogen1}Br
              }
              \qquad
              \chemfig{
                @{la}\ce{FeBr3}
              }
              \arrow{->}
              \chemfig[atom sep=2.5em]{
                [:0]Br
                -[@{bond1}:0]@{halogen2}\pcharge{-90}{Br}
                -[:0]\ncharge{90}{Fe}\ce{Br3}
              }
              \arrow{->}
              \chemfig{
                [:0]\pcharge{45}{Br}
              }
              \+
              \chemfig[atom sep=2.5em]{
                [:0]Br
                -[:0]\ncharge{90}{\ce{FeBr3}}
              }
            }
            \chemmovearrow{
              \reactionfigarrow{halogen1}{-60}{la}{-120}{Red}
              \reactionfigarrow{bond1}{100}{halogen2}{80}{Red}
            }
          }{
            \scheme{
              \chemfig[atom sep=2em]{
                *6(=-=-=-)
              }
              \+
              \chemfig{
                [:0]\ce{Br2}
              }
              \qquad \qquad
              \arrow{->[cat \ce{FeBr3}][
                \shortstack{
                  electrophilic \\
                  \ce{Br+} attack
                }
              ]}
              \qquad
              \chemfig[atom sep=2em]{
                *6(
                (\pcharge{-150}{})
                -=-=-
                (-[:180]Br)
                (-[@{bond3}:120]@{hydrogen}H)
                -[@{bond4}:])
              }
              \+
              \chemfig{
                [:0]@{base}Br
                -[@{bond5}:0]\ncharge{90}{Fe}\ce{Br3}
              }
              \arrow{->}
              \chemfig[atom sep=2em]{
                *6(-=-=-
                (-[:150]Br)
                =)
              }
              \+
              \chemfig{
                [:0]HBr
              }
            }
            \chemmovearrow{
              \reactionfigarrow{base}{135}{hydrogen}{45}{Red}
              \reactionfigarrow{bond3}{30}{bond4}{-15}{Red}
              \reactionfigarrow{bond5}{80}{base}{100}{Red}
            }
          }{Bromination of benzene}{fig:ph73a}{0.77}

        \end{enumerate}

      \end{core}

      \begin{core}[label=core:eas_nitration]{Nitration}

        \textbf{Nitration} involves the attachment of a nitro group (\ce{NO2}) to an aromatic ring and serves primarily as a synthetic intermediate step in practice. Consider first the generation of the \ce{NO2+} electrophile, formed upon exposure of nitric acid to sulfuric acid.

        \reactionfig{
          \scheme{
            \chemfig[atom sep=3em]{
              [:0]H@{oxygen1}O
              -[:0]\pcharge{0}{N}
              (=[:45]O)
              (-[:-45]\ncharge{0}{O})
            }
            \+
            \chemfig{
              [:0]@{hydrogen1}H
              -[@{bond1}:0]@{oxygen2}O@{placeholder}\ce{SO3H}
            }
            \arrow{->}
            \chemfig[atom sep=3em]{
              [:0]\ce{H2}@{oxygen3}\pcharge{90}{O}
              -[@{bond2}:0]\pcharge{0}{N}
              (=[:45]O)
              (-[@{bond3}:-45]@{oxygen4}\ncharge{-45}{O})
            }
            \+
            \chemfig{
              [:0]\ncharge{90}{O}\ce{SO3H}
            }
            \arrow{->}
            \chemfig[atom sep=2em]{
              [:0]O
              =[:0]\pcharge{90}{N}
              =[:0]O
            }
            \+
            \chemfig{
              [:0]\ce{H2O}
            }
            \+
            \chemfig{
              [:0]\ncharge{135}{O}\ce{SO3H}
            }
          }
          \chemmovearrow{
            \reactionfigarrow{oxygen1}{60}{hydrogen1}{120}{Red}
            \reactionfigarrow{bond1}{-100}{oxygen2}{-90}{Red}
            \reactionfigarrow{bond2}{-80}{oxygen3}{-100}{Red}
            \reactionfigarrow{oxygen4}{15}{bond3}{45}{Red}
          }
        }{Generation of the \ce{NO2+} electrophile}{fig:ph74a}{0.77}
        \vspace{-0.5cm}

        With this electrophile established, the nitration of benzene proceeds according to the general mechanism outlined in \autoref{core:eas_general}.

        \reactionfig{
          \scheme{
            \chemfig[atom sep=2em]{
              [:0]@{oxygen1}O
              =[@{bond2}:0]@{nitrogen1}\pcharge{90}{N}
              =[:0]O
            }
            \qquad
            \chemfig[atom sep=2em]{
              *6(
              -=-=
              -@{carbon1}
              =[@{bond1}:]
              )
            }
            \qquad
            \arrow{->}
            \qquad
            \chemfig[atom sep=2em]{
              *6(
              (\pcharge{-150}{})
              -=-=-
              (
              -[:180]\pcharge{180}{N}
              (=[:120]O)
              (-[:240]\ncharge{180}{O})
              )
              (-[@{bond3}:120]@{hydrogen}H)
              -[@{bond4}:])
            }
            \+
            \chemfig{
              [:0]@{base}\ce{H2O}
            }
            \arrow{->}
            \qquad
            \chemfig[atom sep=2em]{
              *6(-=-=-
              (
              -[:150]\pcharge{-90}{N}
              (=[:90]O)
              (-[:210]\ncharge{180}{O})
              )
              -)
            }
            \raisebox{0.7cm}{
              \+
            }
            \chemfig{
              [:0]\ce{H3O+}
              -[:90,,,,White]\ce{^-OSO3H}
            }
          }
          \chemmovearrow{
            \reactionfigarrow{bond1}{-10}{carbon1}{0}{Red}
            \reactionfigarrow{carbon1}{150}{nitrogen1}{60}{Red}
            \reactionfigarrow{bond2}{-80}{oxygen1}{-100}{Red}
            \reactionfigarrow{base}{135}{hydrogen}{45}{Red}
            \reactionfigarrow{bond3}{30}{bond4}{-15}{Red}
          }
        }{Nitration of benzene}{fig:ph75a}{0.77}

      \end{core}

      \begin{core}{Sulfonation}

        \textbf{Sulfonation} entails the attachment of a sulfonic acid group (\ce{SO3H}) to an aromatic ring. The initial step involves a reaction between two molecules of \ce{H2SO4}, arising from the use of \textbf{fuming \ce{H2SO4}} as one of the reagents.

        \reactionfig{
          \scheme{
            \chemfig[atom sep=2em]{
              [:0]H
              -[:0]OS
              (=[:90]O)
              (=[:-90]O)
              -[:0]@{oxygen1}O
              -[:0]H
            }
            \+
            \chemfig{
              [:0]@{hydrogen1}H
              -[@{bond1}:0]@{oxygen2}O\ce{SO3H}
            }
            \arrow{->}
            \chemfig[atom sep=2em]{
              [:0]H@{oxygen3}O
              -[@{bond2}:0]S
              (=[:90]O)
              (=[:-90]O)
              -[@{bond3}:0]@{oxygen4}\pcharge{-90}{O}\ce{H2}
            }
            \+
            \chemfig{
              [:0]\ncharge{90}{O}\ce{SO3H}
            }
            \arrow{->}
            \chemfig[atom sep=2em]{
              [:0]S
              (=[:90]O)
              (=[:-30]O)
              (=[:-150]H\pcharge{-45}{O})
            }
            \+
            \chemfig{
              [:0]\ce{H2O}
            }
            \+
            \chemfig{
              [:0]\ncharge{135}{O}\ce{SO3H}
            }
          }
          \chemmovearrow{
            \reactionfigarrow{oxygen1}{60}{hydrogen1}{120}{Red}
            \reactionfigarrow{bond1}{-100}{oxygen2}{-80}{Red}
            \reactionfigarrow{oxygen3}{-100}{bond2}{-80}{Red}
            \reactionfigarrow{bond3}{100}{oxygen4}{80}{Red}
          }
        }{Formation of fuming \ce{H2SO4}}{fig:ph76a}{0.77}
        \vspace{-0.5cm}

        The reaction then proceeds according to the same general EAS mechanism described in \autoref{core:eas_general}.

        \reactionfig{
          \scheme{
            \chemfig[atom sep=2em]{
              *6(=-=-=-)
            }
            \arrow{->[
              \shortstack{
                fuming \\
                \ce{H2SO4}
              }
            ]}
            \chemfig[atom sep=2em]{
              *6(=-=
              (
              -[:30]S
              (=[:120]O)
              (=[:-60]O)
              -[:30]OH
              )
              -=-)
            }
            \+
            \chemfig{
              [:0]\ce{H3O+}
            }
            \+
            \chemfig{
              [:0] \ce{^-OSO3H}
            }
          }
        }{Sulfonation of benzene}{fig:ph77a}{0.77}

      \end{core}

      \begin{core}[label=core:alkylation]{Alkylation}

        \textbf{Alkylation} refers to the addition of an alkyl group to an aromatic structure. Before an alkane can act as an electrophile, it must first undergo conversion to a carbocation intermediate, as shown in the reactions below.

        \begin{enumerate}

          \item \textbf{Generation of an Electrophile via a Lewis Acid}

          \reactionfig{
            \scheme{
              \chemfig[atom sep=2em]{
                [:0]
                -[:0]
                (-[:120])
                (-[:-120])
                -[:0]@{cl}Cl
              }
              \+
              \chemfig{
                [:0]@{al}Al\ce{Cl3}
              }
              \arrow{->}
              \chemfig[atom sep=2em]{
                [:0]
                -[:0]
                (-[:120])
                (-[:-120])
                -[@{bond}:0]@{cl2}\pcharge{-90}{Cl}
                -[:0,1.5]\ncharge{-90}{Al}\ce{Cl3}
              }
              \arrow{->}
              \chemfig[atom sep=2em]{
                [:0]\pcharge{-90}{}
                (-[:90])
                (-[:-30])
                (-[:-150])
              }
              \+
              \chemfig[atom sep=2em]{
                [:0]Cl
                -[:0,1.5]\ncharge{-90}{Al}\ce{Cl3}
              }
            }
            \chemmovearrow{
              \reactionfigarrow{cl}{60}{al}{120}{Red}
              \reactionfigarrow{bond}{100}{cl2}{80}{Red}
            }
          }{Mechanism for the formation of a tert-butyl carbocation from tert-butyl chloride using aluminum chloride (\ce{AlCl3}) as a Lewis acid catalyst}{fig:ph78a}{0.95}

          This transformation parallels the initial step of the chlorination pathway presented in \autoref{core:eas_halogenation}.

          \item \textbf{Lewis Acid-Catalyzed Epoxide Ring Opening}

          \reactionfig{
            \scheme{
              \chemfig[atom sep=2.5em]{
                [:-30]*3(-
                (-[:-15])
                (-[:-45])
                -@{oxygen1}O-)
              }
              \+
              \chemfig{
                [:0]@{boron}B\ce{F3}
              }
              \arrow{->}
              \chemfig[atom sep=2.5em]{
                [:-30]*3(-
                (-[:-15])
                (-[:-45])
                -[@{bond}:]@{oxygen2}\pcharge{120}{O}
                (-[:60]\ncharge{135}{B}\ce{F3})
                -)
              }
              \arrow{->}
              \chemfig[atom sep=2.5em]{
                [:0]\ce{F3}\ncharge{90}{B}
                -[:-30]O
                -[:30]
                -[:-30]\pcharge{90}{}
                (-[:30])
                (-[:-90])
              }
            }
            \chemmovearrow{
              \reactionfigarrow{oxygen1}{30}{boron}{150}{Red}
              \reactionfigarrow{bond}{0}{oxygen2}{15}{Red}
            }
          }{Boron trifluoride (\ce{BF3}) coordination to an epoxide, promoting regioselective ring cleavage to generate the more thermodynamically stable tertiary carbocation}{fig:ph79a}{0.95}

          \item \textbf{Acid-Catalyzed Carbocation Formation from an Alcohol}

          \reactionfig{
            \scheme{
              \chemfig[atom sep=2em]{
                [:0]
                -[:30]
                (-[:90])
                =[:-30]
                -[:30]
                (-[:120])
                (-[:60]@{oxygen1}OH)
                -[:-30]
              }
              \+
              \chemfig{
                [:0]@{hydrogen}H
                -[@{bond}:0]@{oxygen2}O\ce{SO3H}
              }
              \arrow{->[][\ce{-^-OSO3H}]}
              \chemfig[atom sep=2em]{
                [:0]
                -[:30]
                (-[:90])
                =[:-30]
                -[:30]
                (-[:120])
                (-[@{bond2}:60]@{oxygen3}\pcharge{90}{O}\ce{H2})
                -[:-30]
              }
              \arrow{->}
              \chemfig{
                [:0]
                -[:30]
                (-[:90])
                =[:-30]
                -[:30]\pcharge{-90}{}
                (-[:90])
                -[:-30]
              }
              \+
              \chemfig{
                [:0]\ce{H2O}
              }
            }
            \chemmovearrow{
              \reactionfigarrow{oxygen1}{60}{hydrogen}{120}{Red}
              \reactionfigarrow{bond}{-100}{oxygen2}{-80}{Red}
              \reactionfigarrow{bond2}{150}{oxygen3}{120}{Red}
            }
          }{Protonation of a tertiary alcohol by sulfuric acid, followed by the expulsion of water to form a tertiary carbocation intermediate}{fig:ph80a}{0.95}

          \item \textbf{Protonation of an Aromatic Ring}

          \reactionfig{
            \scheme{
              \chemfig[atom sep=2em]{
                *6(-@{carbon1}
                =[@{bond1}:]
                -
                =[@{bond2}:]@{carbon2}
                (-[:90])
                --)
              }
              \qquad
              \+
              \qquad
              \chemfig{
                [:0]@{hydrogen}H
                -[@{bond}:0]@{oxygen}O\ce{SO3H}
              }
              \arrow{->}
              \chemfig[atom sep=2em]{
                *6(-=
                -\pcharge{30}{}
                -
                (-[:120])
                (-[:60]H)
                --)
              }
              \qquad
              \+
              \qquad
              \chemfig[atom sep=2em]{
                *6(-
                (-[:-90]H)
                -\pcharge{-30}{}
                -=
                (-[:90])
                --)
              }
            }
            \chemmovearrow{
              \reactionfigarrow{bond1}{80}{carbon1}{100}{Orange}
              \reactionfigarrow{carbon1}{-90}{hydrogen}{-120}{Orange}
              \reactionfigarrow{bond2}{-80}{carbon2}{-100}{Blue}
              \reactionfigarrow{carbon2}{60}{hydrogen}{120}{Blue}
              \reactionfigarrow{bond}{100}{oxygen}{80}{Red}
            }
          }{Electrophilic attack on toluene by\ce{H+}; \textit{the first hydrogenation step follows either the blue or orange \ce{e^-} movement}}{fig:ph81a}{0.95}

        \end{enumerate}

        In certain cases, carbocation rearrangement occurs when the initially formed carbocation is unstable.

        \reactionfig{
          \scheme{
            \chemfig[atom sep=2em]{
              *4(-
              (
              -[:-30]
              -[:30]@{oxygen1}OH
              )
              ---)
            }
            \+
            \chemfig{
              [:0]@{hydrogen}H
              -[@{bond1}:0]@{oxygen2}O\ce{SO3H}
            }
            \arrow(aa--bb){->}
            \chemfig[atom sep=2em]{
              *5(-
              -\pcharge{0}{}
              ---)
            }
            \+
            \chemfig{
              [:0]\ce{H2O}
            }
            \+
            \chemfig{
              [:0]\ncharge{90}{O}\ce{SO3H}
            }
            \arrow(@aa--cc)[-60]
            \chemfig{
              *4(-
              (
              -[:-30]@{carbon2}
              -[@{bond3}:30]@{oxygen3}\pcharge{90}{O}\ce{H2}
              )
              -[@{bond2}:]@{carbon1}
              --)
            }
            \arrow(@cc--@bb){->}
          }
          \chemmovearrow{
            \reactionfigarrow{oxygen1}{60}{hydrogen}{100}{Red}
            \reactionfigarrow{bond1}{-100}{oxygen2}{-80}{Red}
            \reactionfigarrow{bond2}{190}{carbon1}{150}{Red}
            \reactionfigarrow{carbon1}{30}{carbon2}{90}{Red}
            \reactionfigarrow{bond3}{90}{oxygen3}{80}{Red}
          }
        }{Acid-catalyzed dehydration and subsequent ring expansion to generate a more thermodynamically stable cyclopentyl carbocation}{fig:ph82a}{0.85}
        \vspace{-0.4cm}

        These rearranged carbocations subsequently undergo EAS with aromatic structures.

        \vspace{0.4cm}
        \reactionfig{
          \scheme{
            \chemfig[atom sep=2em]{
              *6(=-=-=-)
            }
            \+
            \chemfig[atom sep=2em]{
              *4(-
              (
              -[:-30]
              -[:30]OH
              )
              ---)
            }
            \arrow{->[\ce{H2SO4}]}
            \chemfig[atom sep=2em]{
              [:30]*6(=-
              (-[@{bond1}:60]@{hydrogen}H)
              (
              -[:0]*5(-----)
              )
              -[@{bond2}:]\pcharge{60}{}
              -=-)
            }
            \+
            \chemfig[atom sep=2em]{
              [:0]@{oxygen1}\ncharge{-90}{O}\ce{SO3H}
            }
            \arrow{->}
            \chemfig[atom sep=2em]{
              [:30]*6(=-
              (
              -[:0]*5(-----)
              )
              =
              -=-)
            }
          }
          \chemmovearrow{
            \reactionfigarrow{oxygen1}{150}{hydrogen}{45}{Red}
            \reactionfigarrow{bond1}{-45}{bond2}{-135}{Red}
          }
        }{Electrophilic aromatic substitution of benzene by a rearranged cyclopentyl carbocation}{fig:ph83a}{0.8}
        \vspace{-0.4cm}

        In the case above, illustrating rearrangement from a cyclobutyl to a cyclopentyl carbocation, such rearrangement can be avoided through a stepwise pathway involving \textbf{acylation} (\textit{\autoref{core:acylation}}) followed by reduction with \ce{Zn(Hg)} and \ce{HCl}.

        \begin{enumerate}

          \item \textbf{Acylation}

          \reactionfig{
            \scheme{
              \chemfig[atom sep=2em]{
                *6(=-=-=-)
              }
              \+
              \chemfig[atom sep=2em]{
                [:0]Cl
                -[:30]
                (=[:90]O)
                -[:-30]*4(----)
              }
              \arrow{->[cat \ce{BCl3}]}
              \chemfig[atom sep=2em]{
                [:0]Ph
                -[:30]
                (=[:90]O)
                -[:-30]*4(----)
              }
            }
          }{Acylation}{fig:ph84a}{0.77}
          \vspace{-0.5cm}

          \item \textbf{Reduction}

          \reactionfig{
            \scheme{
              \chemfig[atom sep=2em]{
                [:0]Ph
                -[:30]
                (=[:90]O)
                -[:-30]*4(----)
              }
              \arrow{->[Zn(Hg), HCl]}
              \chemfig[atom sep=2em]{
                [:0]Ph
                -[:30]
                (-[:60,,,,Orange]\textcolor{Orange}{H})
                (-[:120,,,,Orange]\textcolor{Orange}{H})
                -[:-30]*4(----)
              }
            }
          }{Reduction of ketone with Zn(Hg) and HCl}{fig:ph85a}{0.77}
          \vspace{-0.5cm}

        \end{enumerate}

      \end{core}

      \begin{core}[label=core:acylation]{Acylation}

        \textbf{Acylation} refers to the addition of an acyl halide to an aromatic structure. The acyl halide must first be rendered electrophilic through activation by a Lewis acid, as shown in the general reaction outline below.

        \reactionfig{
          \scheme{
            \chemname{
              \raisebox{-1cm}{
                \chemfig{
                  [:0]R
                  -[:30]
                  (=[:90]O)
                  -[:-30]@{halogen1}X
                }
              }
            }{Acyl halide}
            \raisebox{-1cm}{
              \+
              \chemfig{
                [:0]@{leaving}LA
              }
            }
            \arrow{->}
            \chemfig{
              [:0]R
              -[:30]
              (=[@{bond1}:90]@{oxygen}O)
              -[@{bond2}:-30]@{halogen2}\pcharge{90}{X}
              -[:0]\ncharge{90}{LA}
            }
            \arrow{->}
            \chemfig{
              [:0]R
              -[:0]C
              ~[:0]\pcharge{45}{O}
            }
            \+
            \chemfig{
              [:0]X
              -[:0]\ncharge{45}{LA}
            }
          }
          \chemmovearrow{
            \reactionfigarrow{halogen1}{60}{leaving}{120}{Red}
            \reactionfigarrow{oxygen}{150}{bond1}{180}{Red}
            \reactionfigarrow{bond2}{-120}{halogen2}{-100}{Red}
          }
        }{Conversion of an acyl halide to an electrophile}{fig:ph86a}{0.77}

        The resulting \ce{RCO+} species can then undergo EAS with an aromatic ring, as demonstrated below using representative compounds.

        \doublereactionfig{
          \scheme{
            \chemfig[atom sep=2em]{
              [:0]*6(=-=-=-)
            }
            \+
            \chemfig{
              [:0]Ph
              -[:30]
              (=[:90]O)
              -[:-30]Cl
            }
            \arrow{->[cat][\ce{FeCl3}]}
            \chemfig[atom sep=2em]{
              [:0]*6(=-
              =[@{bond1}:]@{carbon1}
              -=-)
            }
            \+
            \chemfig{
              [:0]Ph
              -[:0]@{carbon2}C
              ~[@{bond2}:0]@{oxygen1}\pcharge{90}{O}
            }
            \arrow{->}
          }
          \chemmovearrow{
            \reactionfigarrow{bond1}{-170}{carbon1}{170}{Red}
            \reactionfigarrow{carbon1}{45}{carbon2}{110}{Red}
            \reactionfigarrow{bond2}{-100}{oxygen1}{-80}{Red}
          }
        }{
          \scheme{
            \chemfig[atom sep=2em]{
              *6(=
              -\pcharge{-30}{}
              -[@{bond3}:]
              (-[@{bond2}:90]@{hydrogen}H)
              (
              -[:30]
              (=[:90]O)
              -[:-30]Ph
              )
              -=-)
            }
            \+
            \chemfig{
              [:0]@{cl}Cl
              -[@{bond1}:0]@{iron}\ncharge{90}{Fe}\ce{Cl3}
            }
            \arrow{->[][-\ce{FeCl3}]}
            \chemfig[atom sep=2em]{
              *6(=-=
              (
              -[:30]
              (=[:90]O)
              -[:-30]Ph
              )
              -=-)
            }
            \+
            \chemfig{
              [:0]\ce{HCl}
            }
          }
          \chemmovearrow{
            \reactionfigarrow{bond1}{-100}{iron}{-80}{Red}
            \reactionfigarrow{cl}{100}{hydrogen}{80}{Red}
            \reactionfigarrow{bond2}{170}{bond3}{-170}{Red}
          }
        }{Acylation of benzene}{fig:ph87a}{0.8}

      \end{core}


